\section{Introduction}
\label{sec:intro}

Many systems give one person the authority to control devices at a distance. An operator in
the control center of a power company opens and closes breakers in substations. In this
paper, \emph{command} is this authority, and the \emph{commander} is the person who has it.
The devices get the orders through a network, and networks have failures. Thus the devices
must continue to obey a commander who is isolated from the other parties.

This property has a cost. A device that obeys an isolated commander cannot know the cause
of the outage. A link failure and an attack give the device the same messages. In the
attack, an attacker gets the key of the commander and then cuts the links to the other
parties. The device continues to obey, and the other parties cannot tell it to stop.

The attack on three Ukrainian power companies in December 2015 had these two
parts~\cite{eisac2016ukraine}. The attackers used the consoles of the operators to open
breakers in 27 or more substations, and approximately 225{,}000 customers had no
power~\cite[p.~8]{eisac2016ukraine}. The attackers also cut the remote control. The report
records that ``even if the operator workstations were recovered, remote commands could not
be issued to bring the substations back online''~\cite[p.~9]{eisac2016ukraine}. We use this
attack only as an example. The report does not tell if the orders had signatures.

A second attack uses the succession rule. Many systems give command to a deputy when no
message comes from the commander for a specified time. An attacker who controls the links
can stop the messages of a loyal commander until this timeout occurs. Then the deputy gets
command. If the attacker knows the key of the deputy, the attacker has command. In this
paper, the name of this attack is the \emph{timeout attack}.

\paragraph{Problem}
How does command stop, and how does it go to a different person, when the commander can be
hostile and the attacker controls all links? A solution must have three properties. First,
a loyal commander continues through an outage. Second, a hostile commander does not keep
command after the other officers stop their consent. Third, a timeout does not give command
to a hostile deputy.

The first result of this paper shows that a protocol cannot fully have
the first two properties together. It also gives the relation between them.

\paragraph{Known methods}
Byzantine agreement started with a commander who ``need not be
loyal''~\cite[p.~384]{lamport1982byzantine}. But its protocols give a leader no authority
without a quorum~\cite{castro1999pbft}, and a commander must make decisions without a quorum
during an outage. Control systems for the power grid tolerate hostile servers, but not a
hostile operator~\cite{babay2018spire}. Command authentication examines one order at a
time~\cite{mashima2019delaying}. Leases and credentials with a short life give a maximum for
the time that authority stays after revocation~\cite{gray1989leases,zhou2026hrcred}. We
found no work that gives a minimum, and no succession rule that is safe against the timeout
attack (\cref{sec:related}).

\paragraph{Method}
The work has four steps. First, we did a search of the literature with 15 queries and
examined 300 records (\cref{sec:related}). Second, we made a model in which the attacker
controls all links and knows the keys of some officers (\cref{sec:model}). Third, we
compared runs that give a device the same messages, to find limits that apply to all
protocols (\cref{sec:limits}). Fourth, we made a protocol for this model, wrote proofs of
its properties, and calculated its cost (\cref{sec:protocol,sec:analysis,sec:cost}).

\paragraph{Results}
The paper has four results.
\begin{enumerate}
  \item \emph{A limit for the commander} (\cref{thm:limit}). For each protocol, the tenure
  of a hostile commander at a device is not less than the outage tolerance at that device.
  The result applies to limits on time and to limits on the number of orders.
  \item \emph{A limit for the deputy} (\cref{thm:deputy}). The period that a deputy gets on
  a timeout is also the period that an attacker with the key of the deputy gets.
  \item \emph{A protocol} (\cref{sec:protocol}). In \sys{}, a quorum of officers signs each
  renewal of command, and each device compares the end time with its clock. No message to a
  device is necessary to remove a commander. We show unity of command, a maximum for the
  tenure, and a minimum for the outage tolerance
  (\cref{thm:unity,thm:tenure,thm:tolerance}).
  \item \emph{The cost} (\cref{sec:cost}). We give the time without command and the number
  of signatures as formulas, and we compare \sys{} with four known designs.
\end{enumerate}

\paragraph{When to use the results}
The results apply when three conditions are correct. First, one person at a time controls
the devices. Second, the links between that person and the other officers can
have long outages. Third, a hostile commander can cause damage that is not easy to repair.
If outages are short, a quorum for each order is safer. If no officer can be hostile, a
lease from one server is sufficient (\cref{sec:select}).

\paragraph{Work that is done and work that is not done}
\Cref{tab:status} gives the method that shows each result, and the work that is not done.
The proofs are in this paper, but a second person did not examine them at this time. No
result of this paper comes from a test on a power grid.

\begin{table}[t]
\caption{Results and the method that shows each result. The last column tells if the work
is done at this time.}
\label{tab:status}
\centering
\small
\begin{tabularx}{\columnwidth}{@{}>{\raggedright\arraybackslash}X>{\raggedright\arraybackslash}p{0.30\columnwidth}l@{}}
\toprule
Result & Method & Done \\
\midrule
Limits (\cref{thm:limit,thm:deputy}) & Proof & Yes \\
Properties of \sys{} (\cref{thm:unity,thm:tenure,thm:tolerance}, \cref{prp:fallback}) & Proof & Yes \\
Rules of \sys{} & Random tests (\cref{sec:random}) & Yes \\
Rules of \sys{} & Model check & No \\
Cost of \sys{} & Formulas & Yes \\
Cost of expiry and of timeouts with no attacker & Simulation, previous model & Yes \\
Cost with a hostile commander & Simulation & No \\
Local control (rule D3) & Check by a power engineer & No \\
Related work & Abstracts of 300 records & Yes \\
Related work & Full text of each source & No \\
\bottomrule
\end{tabularx}
\end{table}

The paper has these parts. \Cref{sec:background} gives the background, and
\cref{sec:related} the related work and the open problems. \Cref{sec:model} gives the model,
and \cref{sec:limits} the limits. \Cref{sec:protocol,sec:analysis} give the protocol and
its analysis. \Cref{sec:cost} gives the cost and the tests, and \cref{sec:discussion} the
discussion.
